\section{Intermolecular Forces, Solids, and Liquids}

\subsection{Intermolecular Forces}

\subsubsection{Van der Waals interactions}
\[
\text{dispersion and dipole interactions: typically }<10\,\mathrm{kJ\,mol^{-1}}
\]
Van der Waals forces act between permanent or induced dipoles. They dominate noble gas liquids and nonpolar molecular substances.

\subsubsection{Hydrogen bonding}
\[
\mathrm{X{-}H\cdots Y}, \qquad X,Y\in\{\ce{F},\ce{O},\ce{N}\}
\]
Hydrogen bonding requires hydrogen bound to \(\ce{F}\), \(\ce{O}\), or \(\ce{N}\) interacting with \(\ce{F}\), \(\ce{O}\), or \(\ce{N}\). It raises boiling points and stabilizes structures such as ice and biomolecules.

\subsubsection{Like dissolves like}
\[
\text{polar dissolves polar}, \qquad \text{nonpolar dissolves nonpolar}
\]
Solubility is favored when solute-solvent interactions can replace solute-solute and solvent-solvent interactions.

\subsection{Crystals and Unit Cells}

\subsubsection{Crystal lattice and unit cell}
\[
\text{crystal}=\text{periodic translations of a unit cell}
\]
A unit cell is the smallest repeating volume used to describe a crystalline lattice.

\subsubsection{Atoms per cubic unit cell}
\[
\begin{array}{c c c}
\toprule
\text{cell type} & \text{atoms per cell} & \text{packing fraction}\\
\midrule
\text{simple cubic} & 1 & 0.52\\
\text{body-centered cubic} & 2 & 0.68\\
\text{face-centered cubic} & 4 & 0.74\\
\bottomrule
\end{array}
\]
Corner atoms contribute \(1/8\), face atoms \(1/2\), and body-center atoms \(1\) to one unit cell.

\subsubsection{Density of a unit cell}
\[
\rho=\frac{Z M}{N_{\mathrm A}a^3}
\]
\(Z\) is atoms or formula units per unit cell and \(a\) is the cubic cell edge length. Use consistent units for \(a^3\).

\subsubsection{Cubic cell radius relations}
\[
\begin{array}{c c}
\toprule
\text{cell type} & \text{radius relation}\\
\midrule
\text{simple cubic} & a=2r\\
\text{body-centered cubic} & \sqrt{3}a=4r\\
\text{face-centered cubic} & \sqrt{2}a=4r\\
\bottomrule
\end{array}
\]
These relations connect atomic radius and cell edge length for packed spheres.

\subsection{Phase Changes and Vapor Pressure}

\subsubsection{Phase transition enthalpies}
\[
\Delta H_{\mathrm{sub}}=\Delta H_{\mathrm{fus}}+\Delta H_{\mathrm{vap}}
\]
Sublimation can be viewed as fusion followed by vaporization. Fusion and vaporization enthalpies are positive for the forward phase changes.

\subsubsection{Clausius-Clapeyron equation}
\[
\ln P=-\frac{\Delta H_{\mathrm{vap}}}{RT}+C
\]
Vapor pressure increases with temperature. A plot of \(\ln P\) versus \(1/T\) has slope \(-\Delta H_{\mathrm{vap}}/R\).

\subsubsection{Two-point Clausius-Clapeyron equation}
\[
\ln\left(\frac{P_2}{P_1}\right)
=-\frac{\Delta H_{\mathrm{vap}}}{R}\left(\frac{1}{T_2}-\frac{1}{T_1}\right)
\]
This form eliminates the constant \(C\) and is used to calculate vapor pressure or boiling point.

\subsection{Colligative Properties}

\subsubsection{Raoult's law}
\[
P_1=X_1P_1^\circ=(1-X_2)P_1^\circ
\]
A nonvolatile solute lowers the solvent vapor pressure. \(X_1\) is the solvent mole fraction.

\subsubsection{Molality}
\[
m=\frac{n_{\mathrm{solute}}}{m_{\mathrm{solvent}}(\mathrm{kg})}
\]
Molality is independent of temperature because it uses solvent mass, not solution volume.

\subsubsection{Freezing point depression and boiling point elevation}
\[
\Delta T_f=iK_fm, \qquad \Delta T_b=iK_bm
\]
\(K_f\) and \(K_b\) are solvent constants. The van't Hoff factor \(i\) accounts for dissolved particles when dissociation is relevant.
